\cotefile{15}
\scanpage{195}
2°) $X_{\deleted{\unsure{réd}}}$ domine $\operatorname{Spec}\mathbf Z$, i.e. $X$ contient des pts de car résiduelle nulle.

\deleted{Alors \unsure{en plus de ces} \inserted{\unsure{des}\deleted{\unsure{deux}}} (et \unsure{éventuelles} \unsure{cristallines} \unsure{de} \unsure{car d'inégales car}) on a, pour \unsure{la} \unsure{pt}}

\deleted{\illegible{1} (\unsure{cristallin})}

(i) $\xi_l\in\unsure{\mathbb T}(\mathbf Z_l)$ \quad $\forall\ l\in\mathbf P$, $\xi_l$ provenant d'un \unsure{élément} de $\unsure{T}_\alpha(\unsure{\widetilde{\mathbf Z}_{l\mathbf Z}})$.

$\xi_l$ dépend du choix d'un pt géom quelconque de $X$, et est indép. de ce choix modulo opération du $\pi_1$, notation $\xi_{\bar x,l}\unsure{:=}\xi_{X,l}$.

(ii) Si $X$ a des points $x$ de car. $p>0$, alors on dispose également des $\xi_{x,p}$, $\xi^0_{x,p}$.

Connexion : relation, \deleted{\unsure{de la}} si $y$ est un \unsure{pt|vt} en car. $p>0$ et \deleted{\unsure{de}} un \unsure{pt|vt} de car. $0$, entre $\xi_{y,p}$ et $\xi_{x,p}\otimes_{\mathbf Z_p}W(k)$ [$y$ pt à val. dans corps parfait $k$].

\uline{Ils ne sont en général pas isomorphes}, sans doute. \deleted{\unsure{Et}} \inserted{Si} \deleted{\unsure{est son plus}} $k$ est fini, ils sont cependant \unsure{isomorphes} \unsure{par}\deleted{\illegible{5}} \unsure{devant} \unsure{$k$}.

\margin{\uuline{\deleted{\unsure{Vérifier}}}}

(iii) Si $x$ est un \unsure{pt|vt} à val. dans un corps \uline{de car 0}, on a $\xi_{x,\infty}\in\unsure{T}_\alpha(k)$ par Hodge-de Rham. Il n'y a pas lieu de \inserted{avec ses (2) filtrations} \deleted{\illegible{2}}.

\margin{L'une est celle \unsure{qui|que} donne comme gradués associés les $H^{\unsure{q|i}}(X,\Omega^p)$, l'autre (arithmétique) fait pendant à la filtration canonique des motifs. (Elle est triviale pour les motifs simples) la première filtration est \unsure{originale au} \unsure{foncteur} Hodge-de Rham, \unsure{bien faire attention} qu'elle ne provient pas d'une filtration du motif, c'est une filtration du \unsure{f-}\uline{\unsure{de Rham}}.}

\mspage{194}
